\documentclass[11pt]{article}
\usepackage[a4paper,margin=0.9in]{geometry}
\usepackage{amsmath,amssymb,amsthm,mathtools}
\usepackage[dvipsnames]{xcolor}
\usepackage{enumitem}
\usepackage{booktabs}
\usepackage{longtable}
\usepackage{tcolorbox}
\usepackage{fancyhdr}
\usepackage[colorlinks=true,linkcolor=Blue]{hyperref}
\tcbuselibrary{skins,breakable}

\setlist{itemsep=3pt,topsep=4pt}
\pagestyle{fancy}\fancyhf{}
\lhead{\small Exponential Families --- Crash Course}
\rhead{\small Parametric Inference, B.Stat.\ 3rd yr}
\cfoot{\small\thepage}

\newcommand{\E}{\mathbb{E}}
\newcommand{\Var}{\operatorname{Var}}
\newcommand{\Prob}{\mathbb{P}}
\newcommand{\R}{\mathbb{R}}
\newcommand{\Xbar}{\overline{X}}
\newcommand{\iid}{\overset{\text{iid}}{\sim}}
\newcommand{\pto}{\overset{p}{\longrightarrow}}
\newcommand{\I}{\mathbb{I}}

\newtcolorbox{idea}[1]{colback=Blue!4,colframe=Blue!55,title={#1},
  fonttitle=\bfseries,breakable}
\newtcolorbox{confusion}[1]{colback=BrickRed!5,colframe=BrickRed!60,
  title={Common confusion: #1},fonttitle=\bfseries,breakable}
\newtcolorbox{recipe}[1]{colback=ForestGreen!5,colframe=ForestGreen!55,title={#1},
  fonttitle=\bfseries,breakable}

\newcommand{\prf}{\smallskip\noindent\textbf{Proof.}\ }

\begin{document}

\begin{center}
{\LARGE\bfseries Exponential Families}\\[3pt]
{\large A Crash Course from First Principles}\\[5pt]
{\small Parametric Inference, B.Stat.\ 3rd Year --- Lectures 4--5, and everywhere else}
\end{center}

\vspace{4pt}\hrule\vspace{10pt}

\noindent
Fourth of the crash courses. This one is different in kind: Bayes, testing and minimax are
\emph{topics}, but exponential families are the \emph{structure underneath all of them}. Read
this and a lot of the course stops looking like a list of separate results.

\tableofcontents

%%%%%%%%%%%%%%%%%%%%%%%%%%%%%%%%%%%%%%%%%%%%%%%%%%%%%%%%%
\section{Why this class of distributions exists}

Notice something about the whole course. The same handful of families keeps appearing ---
Binomial, Poisson, Exponential, Normal, Gamma --- and for all of them \emph{everything works}:
\begin{itemize}
\item a low-dimensional sufficient statistic exists,
\item it turns out to be complete, so Lehmann--Scheff\'e applies,
\item the Cram\'er--Rao bound is attained,
\item the family has MLR, so UMP tests exist,
\item there is a tidy conjugate prior with two hyperparameters.
\end{itemize}
Meanwhile two families keep breaking things: the \textbf{Cauchy} (sufficient statistic is the
whole order statistic, not complete, no closed-form MLE, no MLR, no useful conjugate prior)
and $\boldsymbol{U[0,\theta]}$ (Cram\'er--Rao meaningless, non-standard rates).

\begin{idea}{The single fact that explains all of this}
The well-behaved families are all \textbf{exponential families} --- the same object wearing
different clothes. The badly behaved ones are not.

So exponential families are not one more topic. They are the \emph{hypothesis} of most of the
theorems in the course, and the two villains are villains precisely because they fail it.
Everything in Section 5 below is a consequence of one definition.
\end{idea}

%%%%%%%%%%%%%%%%%%%%%%%%%%%%%%%%%%%%%%%%%%%%%%%%%%%%%%%%%
\section{The definition, and how to recognise one}

\begin{idea}{Definition (one parameter)}
$\{f(\cdot\mid\theta):\theta\in\Theta\}$ is a \textbf{one-parameter exponential family} if
\begin{enumerate}[label=(\roman*)]
\item the \textbf{support} $\{x:f(x\mid\theta)>0\}$ does \textbf{not depend on $\theta$}, and
\item $f(x\mid\theta)=p(\theta)\,h(x)\,\exp\{c(\theta)\,T(x)\}$
\end{enumerate}
for some functions $p>0$, $h>0$, $c$, $T$.
\end{idea}

Both conditions matter, and condition (i) is the one people skip.

\begin{recipe}{The recognition drill}
\textbf{Take logs.} You have an exponential family iff
\[
\log f(x\mid\theta)=\underbrace{\log p(\theta)}_{\text{pure }\theta}
+\underbrace{\log h(x)}_{\text{pure }x}
+\underbrace{c(\theta)\,T(x)}_{\text{a product: one }\theta\text{-factor}\times\text{one }x\text{-factor}} ,
\]
with a support that does not move. So: expand $\log f$, sort the terms into three piles, and
check that every term mixing $\theta$ and $x$ separates into a product.

\smallskip
\textbf{Quick test for failure.} If it is an exponential family then
$\dfrac{\partial^{2}}{\partial\theta\,\partial x}\log f=c'(\theta)T'(x)$, which must factor
into (function of $\theta$) $\times$ (function of $x$). If the mixed partial does not
factor, it is not an exponential family.
\end{recipe}

\subsection{Worked identifications}

\textbf{Binomial}, $X\sim\mathrm{Bin}(m,\theta)$:
\[
\binom mx\theta^{x}(1-\theta)^{m-x}
=\underbrace{(1-\theta)^{m}}_{p(\theta)}\underbrace{\binom mx}_{h(x)}
\exp\Big\{\underbrace{\log\tfrac{\theta}{1-\theta}}_{c(\theta)}\underbrace{x}_{T(x)}\Big\}.
\]
Support $\{0,\dots,m\}$, free of $\theta$. \checkmark

\smallskip
\textbf{Poisson}: $\dfrac{e^{-\theta}\theta^{x}}{x!}
=\underbrace{e^{-\theta}}_{p}\underbrace{\tfrac{1}{x!}}_{h}
\exp\{\underbrace{\log\theta}_{c}\cdot\underbrace{x}_{T}\}$.

\smallskip
\textbf{Exponential (mean $\theta$)}: $\dfrac1\theta e^{-x/\theta}
=\underbrace{\theta^{-1}}_{p}\underbrace{\I[x>0]}_{h}\exp\{\underbrace{-\tfrac1\theta}_{c}\cdot
\underbrace{x}_{T}\}$. Support $(0,\infty)$, free of $\theta$. \checkmark

\smallskip
\textbf{$N(\theta,1)$}: $\dfrac{1}{\sqrt{2\pi}}e^{-(x-\theta)^2/2}
=\underbrace{\tfrac{e^{-\theta^2/2}}{\sqrt{2\pi}}}_{p}\underbrace{e^{-x^2/2}}_{h}
\exp\{\underbrace{\theta}_{c}\cdot\underbrace{x}_{T}\}$.

\smallskip
\textbf{Geometric}, $f=\theta(1-\theta)^x$ on $x=0,1,\dots$:
$=\underbrace{\theta}_{p}\cdot\underbrace{1}_{h}\exp\{\underbrace{\log(1-\theta)}_{c}\cdot
\underbrace{x}_{T}\}$.

\smallskip
\textbf{Negative binomial} (the HW1 Q1 form),
$f=\tfrac{\Gamma(\alpha+x)}{\Gamma(\alpha)x!}\theta^{\alpha}(1-\theta)^{x}$:
$p=\theta^{\alpha}$, $h=\tfrac{\Gamma(\alpha+x)}{\Gamma(\alpha)x!}$, $c=\log(1-\theta)$, $T=x$.

\smallskip
\textbf{Power series family}, $f=\dfrac{a_x\theta^{x}}{A(\theta)}$:
$p=1/A(\theta)$, $h=a_x$, $c=\log\theta$, $T=x$. (So Bernoulli, Binomial, Poisson, Geometric
and Negative Binomial are all power-series families, hence all exponential families --- one
observation covering five rows of the table.)

\subsection{And the failures}

\textbf{$U[0,\theta]$.} $f=\tfrac1\theta\I[0\le x\le\theta]$. The functional form is fine, but
the \textbf{support moves with $\theta$}: condition (i) fails outright. Not an exponential
family.

\smallskip
\textbf{$U[\theta,\theta+1]$.} Same reason.

\smallskip
\textbf{Cauchy$(\theta)$.} The support is all of $\R$, so condition (i) holds --- the failure
is structural. Apply the mixed-partial test with $u=x-\theta$:
\[
\log f=-\log\pi-\log(1+u^{2}),\qquad
\frac{\partial}{\partial x}\log f=\frac{-2u}{1+u^{2}},\qquad
\frac{\partial^{2}}{\partial\theta\,\partial x}\log f=\frac{2(1-u^{2})}{(1+u^{2})^{2}} .
\]
This vanishes exactly at $x=\theta\pm1$ --- zeros whose \emph{position moves with $\theta$}.
A product $c'(\theta)T'(x)$ has its $x$-zeros at fixed places. Contradiction, so the Cauchy is
not an exponential family.

\emph{Second proof, and the more memorable one:} in an exponential family
$\sum_iT(X_i)$ is sufficient for every $n$ --- a \textbf{one-dimensional} statistic no matter
how much data you have. For the Cauchy the minimal sufficient statistic is the whole order
statistic, whose dimension grows with $n$. So it cannot be an exponential family.

\begin{confusion}{``exponential distribution'' vs ``exponential family''}
Unrelated names. The \emph{exponential distribution} $\tfrac1\theta e^{-x/\theta}$ is one
member of the exponential \emph{family} class --- but so are the Binomial, Poisson, Normal and
Gamma, none of which is ``an exponential distribution''. The word refers to the
$\exp\{c(\theta)T(x)\}$ \emph{form}, not to any particular density.

Equally: having an $e^{(\cdot)}$ in the formula neither proves nor disproves membership.
$U[0,\theta]$ has no exponential and is not a member; the Cauchy has no exponential and is
also not a member --- but for a completely different reason.
\end{confusion}

%%%%%%%%%%%%%%%%%%%%%%%%%%%%%%%%%%%%%%%%%%%%%%%%%%%%%%%%%
\section{Natural parameter and natural parameter space}

The parametrisation you are handed is arbitrary. The theory is cleanest in the parameter that
appears \emph{linearly in the exponent}.

\begin{idea}{Definitions}
$\eta:=c(\theta)$ is the \textbf{natural parameter}, and
\[
\Gamma:=\{c(\theta):\theta\in\Theta\}\subseteq\R
\]
is the \textbf{natural parameter space}. Rewriting the density in terms of $\eta$ gives the
\emph{natural form}
\[
f(x\mid\eta)=h(x)\exp\{\eta\,T(x)-A(\eta)\},\qquad
A(\eta)=\log\int h(x)e^{\eta T(x)}\,dx ,
\]
where $A$ (the \emph{log-partition} or \emph{cumulant} function) is just whatever makes the
density integrate to 1.
\end{idea}

\begin{confusion}{$\Theta$ versus $\Gamma$}
The completeness criterion is a condition on $\boldsymbol\Gamma$, \textbf{not} on $\Theta$.
$\Theta$ is almost always a nice interval; that tells you nothing. You must compute the image
$c(\Theta)$.
\begin{center}
\renewcommand{\arraystretch}{1.3}
\begin{tabular}{@{}llll@{}}
\toprule
Family & $\Theta$ & $c(\theta)$ & $\Gamma=c(\Theta)$ \\
\midrule
$\mathrm{Bin}(m,\theta)$ & $(0,1)$ & $\log\frac{\theta}{1-\theta}$ & $\R$ \\
$\mathrm{Poi}(\theta)$ & $(0,\infty)$ & $\log\theta$ & $\R$ \\
$\mathrm{Exp}(\theta)$ & $(0,\infty)$ & $-1/\theta$ & $(-\infty,0)$ \\
Power series & $(0,b)$ & $\log\theta$ & $(-\infty,\log b)$ \\
\bottomrule
\end{tabular}
\end{center}
In each case $\Gamma$ contains an open interval, which is all the criterion asks.
\end{confusion}

\subsection{Multiparameter version}

\[
f(x\mid\theta)=p(\theta)h(x)\exp\Big\{\sum_{j=1}^{k}c_j(\theta)T_j(x)\Big\},
\qquad
\Gamma=\big\{(c_1(\theta),\dots,c_k(\theta)):\theta\in\Theta\big\}\subseteq\R^{k}.
\]
For $n$ iid observations the sufficient statistic is
$S=\big(\sum_iT_1(X_i),\dots,\sum_iT_k(X_i)\big)$, and for completeness $\Gamma$ must contain a
\textbf{$k$-dimensional rectangle with non-empty interior}.

\smallskip
\textbf{$N(\mu,\sigma^{2})$:}
\[
f=\underbrace{\frac{e^{-\mu^{2}/2\sigma^{2}}}{\sqrt{2\pi\sigma^{2}}}}_{p(\theta)}\cdot
\underbrace{1}_{h}\cdot
\exp\Big\{\underbrace{\tfrac{\mu}{\sigma^{2}}}_{c_1}\underbrace{x}_{T_1}
+\underbrace{\big(-\tfrac{1}{2\sigma^{2}}\big)}_{c_2}\underbrace{x^{2}}_{T_2}\Big\},
\quad \Gamma=\R\times(-\infty,0).
\]
An open half-plane, so it contains a rectangle: $(\sum X_i,\sum X_i^{2})$ is complete
sufficient.

\smallskip
\textbf{Gamma$(\alpha,\beta)$} (shape, scale, both unknown): $T=(\log x,\ x)$,
$c=(\alpha-1,\ -1/\beta)$, $\Gamma=(-1,\infty)\times(-\infty,0)$. Complete.

\smallskip
\textbf{Beta$(\theta,\phi)$}: $T=(\log x,\ \log(1-x))$, $c=(\theta-1,\ \phi-1)$,
$\Gamma=(-1,\infty)^{2}$. Complete, with
\[
S=\Big(\sum_i\log X_i,\ \sum_i\log(1-X_i)\Big)
\ \equiv\ \Big(\prod_iX_i,\ \prod_i(1-X_i)\Big).
\]
\emph{That is why the class notes record the Beta sufficient statistic as a pair of
products}: sums of logs are logs of products.

%%%%%%%%%%%%%%%%%%%%%%%%%%%%%%%%%%%%%%%%%%%%%%%%%%%%%%%%%
\section{Curved families: where the criterion bites}

\begin{idea}{Full rank versus curved}
Write $k$ for the number of $(c_j,T_j)$ terms and $d$ for the dimension of $\Theta$.
\begin{itemize}
\item If $d=k$ and $\Gamma$ has non-empty interior in $\R^{k}$, the family is \textbf{full
rank} and the criterion applies: $S$ is complete.
\item If $d<k$, the natural parameters are constrained to a lower-dimensional surface. The
family is \textbf{curved}, $\Gamma$ contains no rectangle, and \textbf{the criterion says
nothing}.
\end{itemize}
\end{idea}

\textbf{The standard example: $N(\theta,\theta^{2})$}, $\theta>0$. Here $k=2$ but $d=1$:
\[
(\eta_1,\eta_2)=\Big(\frac{\theta}{\theta^{2}},\ -\frac{1}{2\theta^{2}}\Big)
=\Big(\frac1\theta,\ -\frac{1}{2\theta^{2}}\Big),
\qquad\text{so}\qquad \eta_2=-\tfrac12\eta_1^{2}:
\]
a \emph{parabola} in $\R^{2}$, which contains no rectangle. The criterion is unavailable ---
and in this case completeness genuinely fails.

\begin{idea}{Proof that $(\sum X_i,\sum X_i^{2})$ is not complete for $N(\theta,\theta^{2})$}
Exhibit two \emph{different} unbiased estimators of $\theta^{2}$, both functions of $S$. Using
$\E X_i^{2}=\Var+\text{mean}^2=\theta^{2}+\theta^{2}=2\theta^{2}$ and
$\E\Xbar^{2}=\tfrac{\theta^{2}}{n}+\theta^{2}=\tfrac{(n+1)\theta^{2}}{n}$:
\[
\E\Big[\frac{1}{2n}\sum_iX_i^{2}\Big]=\theta^{2},
\qquad
\E\Big[\frac{n}{n+1}\,\Xbar^{2}\Big]=\theta^{2}.
\]
So $g(S)=\dfrac{1}{2n}\sum_iX_i^{2}-\dfrac{n}{n+1}\Xbar^{2}$ satisfies
$\E_\theta\,g(S)=0$ for every $\theta$, while $g\not\equiv0$. That is a non-trivial zero
function: \textbf{$S$ is not complete}. (Simulation at $\theta=2$, $n=6$: $\E g=-0.001$ against
a standard deviation of $1.61$ --- mean zero, but very far from the zero function.)

Consequence: Lehmann--Scheff\'e is unavailable here, and neither estimator above is a UMVUE.
\end{idea}

%%%%%%%%%%%%%%%%%%%%%%%%%%%%%%%%%%%%%%%%%%%%%%%%%%%%%%%%%
\section{The five consequences --- why the course cares}

This is the payoff. \textbf{One structure, five theorems.}

\begin{idea}{1. Sufficiency of fixed dimension}
For $X_1,\dots,X_n$ iid,
\[
f(\mathbf x\mid\theta)=p(\theta)^{n}\Big(\prod_ih(x_i)\Big)
\exp\Big\{c(\theta)\sum_iT(x_i)\Big\},
\]
so by the factorisation theorem $S=\sum_iT(X_i)$ is sufficient --- \textbf{one number,
whatever $n$ is}. Contrast the Cauchy, where you must keep all $n$ order statistics. Data
reduction of this kind is exactly what exponential families are for.
\end{idea}

\begin{idea}{2. Completeness, hence Lehmann--Scheff\'e}
If $\Gamma$ contains an open interval (rectangle in the $k$-parameter case), $S$ is
\textbf{complete}. Combined with (1), Lehmann--Scheff\'e applies: \emph{any} unbiased function
of $S$ is automatically \emph{the} UMVUE.

That is why the UMVUE recipe is so short for these families --- find one unbiased $h(S)$ by
series matching, and stop.
\end{idea}

\begin{idea}{3. The Cram\'er--Rao bound is attained --- for the right target}
In natural form, $\log f=\log h(x)+\eta T(x)-A(\eta)$, so the score is
\[
\frac{\partial}{\partial\eta}\log f=T(x)-A'(\eta)=T(x)-\E_\eta[T],
\]
which is exactly the factorisation $k(\eta)[T-\psi(\eta)]$ required for equality in
Cauchy--Schwarz. \textbf{So $T$ attains the CRLB for $\psi(\eta)=\E_\eta[T]$}, and it is the
UMVUE by (2).

The caveat that costs marks: this holds for $\E[T]$, \textbf{not for arbitrary $\psi$}. The
UMVUE of $e^{-\lambda}$ in the Poisson (Mock Q3) is a function of $S$ but is \emph{not} $S$
itself, and its variance sits strictly above the bound.
\end{idea}

\begin{idea}{4. MLR, hence UMP tests}
For $\theta_0<\theta_1$,
\[
\frac{f_{\theta_1}(x)}{f_{\theta_0}(x)}
=\frac{p(\theta_1)}{p(\theta_0)}\exp\big\{[c(\theta_1)-c(\theta_0)]T(x)\big\},
\]
increasing in $T(x)$ whenever $c$ is non-decreasing. So the family has \textbf{MLR in $T$},
and Karlin--Rubin delivers a UMP test for every one-sided hypothesis. The Cauchy, not being an
exponential family, has no MLR and no UMP test.
\end{idea}

\begin{idea}{5. Finite-dimensional conjugate priors}
$\prod_if(x_i\mid\theta)$ depends on the data only through $\big(n,\sum_iT(x_i)\big)$ --- two
numbers. So the family of priors
\[
\pi(\theta)\ \propto\ p(\theta)^{a}\exp\{c(\theta)\,b\},\qquad (a,b)\ \text{hyperparameters},
\]
is closed under sampling, with the update $(a,b)\mapsto(a+n,\ b+\sum_iT(x_i))$ ---
\emph{pure arithmetic}. That is Beta--Binomial, Gamma--Poisson and Normal--Normal, all at
once; look back at those update rules and you will see exactly this shape.

For the Cauchy no such compression exists, so the conjugate family must grow by one
hyperparameter per observation --- useless, which is the point of the Lecture 7 discussion.
\end{idea}

\begin{center}
\renewcommand{\arraystretch}{1.4}
\begin{tabular}{@{}p{4.3cm}p{5.3cm}p{5.5cm}@{}}
\toprule
& \textbf{Exponential family} & \textbf{Cauchy / $U[0,\theta]$} \\
\midrule
Sufficient statistic & $\sum_iT(X_i)$, fixed dimension &
order statistic (Cauchy); $X_{(n)}$ ($U[0,\theta]$) \\
Complete? & yes, if $\Gamma$ has interior & no (Cauchy); yes but proved directly ($U[0,\theta]$) \\
CRLB & attained for $\E[T]$ & regularity fails ($U[0,\theta]$) \\
MLR / UMP & yes, if $c$ monotone & no (Cauchy); yes, directly ($U[0,\theta]$) \\
Conjugate prior & two hyperparameters & grows with $n$ (Cauchy); Pareto ($U[0,\theta]$) \\
\bottomrule
\end{tabular}
\end{center}

\begin{confusion}{``So nothing works outside exponential families?''}
Not so --- and $U[0,\theta]$ is the proof. It is \emph{not} an exponential family, yet
$X_{(n)}$ is complete sufficient, the family has MLR, a UMP test exists (even a two-sided
one), and there is a Pareto conjugate prior.

The correct statement is: being an exponential family is a \textbf{sufficient} condition for
these conclusions, not a necessary one. When the hypothesis fails you simply lose the general
theorem and must verify the conclusion \emph{by hand} --- which is precisely what the
$U[0,\theta]$ questions in this course make you do.
\end{confusion}

%%%%%%%%%%%%%%%%%%%%%%%%%%%%%%%%%%%%%%%%%%%%%%%%%%%%%%%%%
\section{The cumulant trick: moments for free}

\begin{idea}{The identity}
In natural form $f(x\mid\eta)=h(x)e^{\eta T(x)-A(\eta)}$,
\[
\boxed{\ \E_\eta[T]=A'(\eta),\qquad \Var_\eta(T)=A''(\eta),\qquad I(\eta)=A''(\eta).\ }
\]
\end{idea}

\prf Integrating the density to 1 gives $\int h(x)e^{\eta T(x)}dx=e^{A(\eta)}$. Differentiate
in $\eta$ and divide by $e^{A(\eta)}$:
\[
\int T(x)h(x)e^{\eta T(x)-A(\eta)}dx=A'(\eta)\ \Longrightarrow\ \E[T]=A'(\eta).
\]
Differentiate once more: $\E[T^{2}]=A''+(A')^{2}$, so $\Var(T)=A''(\eta)$. And since the score
is $T-A'(\eta)$, the Fisher information is $\Var(T)=A''(\eta)$. $\square$

\textbf{Poisson, as a check.} $\eta=\log\lambda$, $T=x$, $h=1/x!$, and
$A(\eta)=\lambda=e^{\eta}$. Then $\E[X]=A'=e^{\eta}=\lambda$ and $\Var(X)=A''=e^{\eta}=\lambda$
--- both classical facts in one line, with no sums.

\begin{idea}{This subsumes the power-series identities of HW1 Q4}
For a power-series family, $\eta=\log\theta$ (so $\tfrac{d\theta}{d\eta}=\theta$) and
$A(\eta)=\log \tilde A(e^{\eta})$ where $\tilde A(\theta)=\sum_xa_x\theta^{x}$. Then
\[
\E[X]=\frac{dA}{d\eta}=\frac{\tilde A'(\theta)}{\tilde A(\theta)}\cdot\theta=\psi(\theta),
\qquad
\Var(X)=\frac{d^{2}A}{d\eta^{2}}=\frac{d\psi}{d\eta}
=\frac{d\psi}{d\theta}\cdot\frac{d\theta}{d\eta}=\theta\,\psi'(\theta).
\]
Those are exactly the two boxed identities from Homework 1 Q4 ---
$\E X=\theta \tilde A'/\tilde A$ and $\Var X=\theta\psi'(\theta)$ --- obtained here with no
computation at all. And the CRLB check falls out too:
$\text{CRLB}=[\psi']^{2}/I=[\psi']^{2}\theta/\psi'=\theta\psi'=\Var(X)$, attained. \checkmark
\end{idea}

%%%%%%%%%%%%%%%%%%%%%%%%%%%%%%%%%%%%%%%%%%%%%%%%%%%%%%%%%
\section{Master table}

\renewcommand{\arraystretch}{1.45}
\begin{longtable}{@{}p{2.9cm}p{2.0cm}p{2.6cm}p{3.0cm}p{3.4cm}@{}}
\toprule
Family & $T(x)$ & $c(\theta)$ & $\Gamma$ & Complete suff.\ stat. \\
\midrule
\endhead
$\mathrm{Bern}/\mathrm{Bin}(m,\theta)$ & $x$ & $\log\frac{\theta}{1-\theta}$ & $\R$ & $\sum X_i$ \\
$\mathrm{Poi}(\theta)$ & $x$ & $\log\theta$ & $\R$ & $\sum X_i$ \\
$\mathrm{Geom}(\theta)$ & $x$ & $\log(1-\theta)$ & $(-\infty,0)$ & $\sum X_i$ \\
$\mathrm{NegBin}(\alpha,\theta)$ & $x$ & $\log(1-\theta)$ & $(-\infty,0)$ & $\sum X_i$ \\
Power series & $x$ & $\log\theta$ & $(-\infty,\log b)$ & $\sum X_i$ \\
$\mathrm{Exp}(\theta)$ & $x$ & $-1/\theta$ & $(-\infty,0)$ & $\sum X_i$ \\
$N(\theta,1)$ & $x$ & $\theta$ & $\R$ & $\Xbar$ \\
$N(\mu,\sigma^{2})$ & $(x,x^{2})$ & $(\frac{\mu}{\sigma^{2}},-\frac{1}{2\sigma^{2}})$ &
$\R\times(-\infty,0)$ & $(\sum X_i,\sum X_i^{2})$ \\
$N(0,\sigma^{2})$ & $x^{2}$ & $-\frac{1}{2\sigma^{2}}$ & $(-\infty,0)$ & $\sum X_i^{2}$ \\
$\mathrm{Gamma}(\alpha,\beta)$ & $(\log x,x)$ & $(\alpha-1,-\frac1\beta)$ & open rectangle &
$(\sum\log X_i,\sum X_i)$ \\
$\mathrm{Beta}(\theta,\phi)$ & $(\log x,\log(1-x))$ & $(\theta-1,\phi-1)$ & $(-1,\infty)^{2}$ &
$(\prod X_i,\prod(1-X_i))$ \\
Multinomial & $(y_1,\dots,y_{k-1})$ & $\log\frac{p_j}{1-p_1-\cdots}$ & $\R^{k-1}$ &
$(Y_1,\dots,Y_{k-1})$ \\
\midrule
\multicolumn{5}{@{}l}{\textbf{Not exponential families}} \\
$N(\theta,\theta^{2})$ & \multicolumn{4}{l}{curved: $\Gamma$ is the parabola
$\eta_2=-\tfrac12\eta_1^{2}$; sufficient but \emph{not} complete} \\
$U[0,\theta]$, $U[\theta,\theta+1]$ & \multicolumn{4}{l}{support depends on $\theta$;
handle by hand} \\
$\mathrm{Cauchy}(\theta)$ & \multicolumn{4}{l}{mixed partial does not factor; sufficient
statistic grows with $n$} \\
\bottomrule
\end{longtable}

%%%%%%%%%%%%%%%%%%%%%%%%%%%%%%%%%%%%%%%%%%%%%%%%%%%%%%%%%
\section{Practice, with answers}

\begin{enumerate}[label=\textbf{P\arabic*.}]
\item $X_1,\dots,X_n\iid\mathrm{Geometric}$, $f(x\mid\theta)=\theta(1-\theta)^{x}$,
$x=0,1,2,\dots$ Show this is an exponential family, identify $\Gamma$, and deduce that
$\sum X_i$ is complete sufficient. Then find the UMVUE of $\theta$ for $n=1$.

\item Show that $N(\theta,\theta^{2})$ is a curved exponential family, and exhibit a
non-trivial zero function to show its sufficient statistic is not complete.

\item $X_1,\dots,X_n\iid\mathrm{Gamma}(\alpha,\beta)$ with both parameters unknown. Identify
$T$, $c$, and $\Gamma$, and state the complete sufficient statistic.

\item Use $\E[T]=A'(\eta)$, $\Var(T)=A''(\eta)$ to obtain the mean and variance of the
Binomial from its cumulant function.

\item $U[0,\theta]$ is not an exponential family. Which conclusions of Section 5 nevertheless
hold for it, and what does that tell you about the logical status of the theory?
\end{enumerate}

\subsection*{Answers}

\noindent\textbf{P1.} $f=\theta\exp\{x\log(1-\theta)\}$, so $p=\theta$, $h=1$, $T=x$,
$c=\log(1-\theta)$. Support $\{0,1,2,\dots\}$ is free of $\theta$. \checkmark\ As $\theta$
ranges over $(0,1)$, $\Gamma=(-\infty,0)$ --- an open interval, so $\sum X_i$ is complete
sufficient.

For $n=1$, unbiasedness reads $\sum_x T(x)\theta(1-\theta)^{x}=\theta$, i.e.\ with $p=1-\theta$,
$\sum_x T(x)p^{x}=1$ for all $p\in(0,1)$. Matching coefficients: $T(0)=1$ and $T(x)=0$ for
$x\ge1$. So the UMVUE is $\boxed{\I\{X=0\}}$, with variance $\theta(1-\theta)$. (Sensible: you
see one observation, and the only unbiased signal about $\theta=\Prob(X=0)$ is whether $X$ was
$0$.)

\smallskip
\noindent\textbf{P2.} With $\mu=\theta$, $\sigma^{2}=\theta^{2}$,
$(\eta_1,\eta_2)=(1/\theta,\ -1/(2\theta^{2}))$, so $\eta_2=-\tfrac12\eta_1^{2}$: the natural
parameters lie on a parabola, a one-dimensional curve in $\R^{2}$. It contains no rectangle,
so the completeness criterion does not apply --- the family is \emph{curved}.

Completeness genuinely fails. Since $\E X_i^{2}=2\theta^{2}$ and
$\E\Xbar^{2}=\tfrac{(n+1)}{n}\theta^{2}$, both
$\tfrac{1}{2n}\sum_iX_i^{2}$ and $\tfrac{n}{n+1}\Xbar^{2}$ are unbiased for $\theta^{2}$, so
\[
g(S)=\frac{1}{2n}\sum_iX_i^{2}-\frac{n}{n+1}\Xbar^{2}
\]
has $\E_\theta g=0$ for all $\theta$ with $g\not\equiv0$. Hence $S$ is not complete.

\smallskip
\noindent\textbf{P3.}
$f=\dfrac{x^{\alpha-1}e^{-x/\beta}}{\Gamma(\alpha)\beta^{\alpha}}
=\dfrac{1}{\Gamma(\alpha)\beta^{\alpha}}\exp\{(\alpha-1)\log x-\tfrac{x}{\beta}\}$, so
$T=(\log x,\ x)$ and $c=(\alpha-1,\ -1/\beta)$, giving
$\Gamma=(-1,\infty)\times(-\infty,0)$ --- an open rectangle. Therefore
$S=\big(\sum_i\log X_i,\ \sum_iX_i\big)$ is complete sufficient.

\smallskip
\noindent\textbf{P4.} $X\sim\mathrm{Bin}(m,\theta)$ has $\eta=\log\frac{\theta}{1-\theta}$, so
$\theta=\frac{e^{\eta}}{1+e^{\eta}}$ and, from $p(\theta)=(1-\theta)^{m}$,
\[
A(\eta)=-\log p=-m\log(1-\theta)=m\log(1+e^{\eta}).
\]
Then
\[
\E[X]=A'(\eta)=\frac{me^{\eta}}{1+e^{\eta}}=m\theta,\qquad
\Var(X)=A''(\eta)=\frac{me^{\eta}}{(1+e^{\eta})^{2}}=m\theta(1-\theta). \ \checkmark
\]

\smallskip
\noindent\textbf{P5.} It fails condition (i) --- the support $[0,\theta]$ moves. Yet:
$X_{(n)}$ is sufficient (factorisation) \emph{and} complete (proved directly by
differentiating $\int_0^\theta g(u)nu^{n-1}du=0$ in $\theta$); the family has MLR in $X_{(n)}$
(the ratio is constant then jumps to $+\infty$); a UMP test exists, and unusually a two-sided
one as well; and the Pareto family is conjugate. What genuinely breaks is Cram\'er--Rao, which
needs the regularity that moving support destroys.

\textbf{Moral:} the exponential-family hypothesis is \emph{sufficient} for the good behaviour,
not \emph{necessary}. Losing it means losing the general theorem, so you must verify the
conclusion by hand --- which is exactly why $U[0,\theta]$ appears so often in this course.

\vfill
\begin{center}
\small\itshape
Companions: \texttt{PI\_Bayes\_Crash\_Course.pdf},
\texttt{PI\_Testing\_Crash\_Course.pdf},
\texttt{PI\_Minimax\_and\_Theorems.pdf};
\texttt{PI\_Exam\_Revision.pdf} \S5--\S7.
\end{center}

\end{document}
